\documentclass[12pt]{article}

% Figura auxiliar para o problema do fluxo através do cone
% Mexilhoeira Grande, Sex Nov 30 20:47:38 WET 2018
% Written by Tordar

\usepackage{graphicx}
\usepackage{pst-all}
\usepackage{pst-slpe}

\pagestyle{empty} 

\begin{document}

\begin{pspicture}
\psset{unit=1cm}
\psline[linewidth=1pt,linestyle=dashed,dash=3mm 1mm](0,2)(2,2)(2,0)
\psline[linewidth=3pt,linecolor=black]{<->}(5,0)(0,0)(0,5)
\psline[linewidth=3pt,linecolor=green]{-}(0,4)(4,0)
\psline[linewidth=4pt,linecolor=blue]{->}(2,2)(3,3)
\psline[linewidth=2pt,linecolor=black]{<->}(0,-0.5)(2,-0.5)
\psline[linewidth=2pt,linecolor=black]{<->}(0,-1.5)(4,-1.5)
\psline[linewidth=2pt,linecolor=black]{<->}(-0.5,0)(-0.5,2)
\psline[linewidth=2pt,linecolor=black]{<->}(-0.5,2)(-0.5,4)
\psline[linewidth=2pt,linecolor=black]{<->}(-2.5,0)(-2.5,4)
\psarc[linewidth=3pt,linecolor=blue]{->}(0,4){1.5}{-90}{-45}
\rput{0}(5.5,0){\scalebox{1.5}{$X$}}
\rput{0}(0,5.5){\scalebox{1.5}{$Y$}}
\rput{0}(3.2,3.2){\scalebox{1.5}{$\mathbf{E}_n$}}
\rput{0}(1,-1){\scalebox{1.5}{$r$}}
\rput{0}(2,-2){\scalebox{1.5}{$R$}}
\rput{0}(-1,1){\scalebox{1.5}{$z$}}
\rput{0}(-1.5,2.5){\scalebox{1.5}{$h-z$}}
\rput{0}(-3,2){\scalebox{1.5}{$h$}}
\rput{0}(0.75,2.4){\scalebox{1.5}{$\alpha$}}
\end{pspicture}

\end{document}
